% =============================================================================
% Project: 302006 - OOD-robustness benchmark of single-cell foundation models
% File: supplementary.tex (P2 rewrite, assembled 2026-09-12 from
%   revision/drafts/SI_resolved_data.tex + tables_v3.md + result CSVs)
% COMPILE ORDER (so cross-references resolve):
%   pdflatex supplementary.tex
%   pdflatex supplementary.tex
%   pdflatex main.tex
%   pdflatex main.tex
% =============================================================================

\documentclass[11pt,a4paper]{article}

\usepackage[utf8]{inputenc}
\usepackage[T1]{fontenc}
\usepackage{amsmath,amssymb,amsfonts}
\usepackage{graphicx}
\usepackage{booktabs}
\usepackage{multirow}
\usepackage{array}
\usepackage{xcolor}
\usepackage{hyperref}
\usepackage[margin=1in]{geometry}
\usepackage{natbib}
\usepackage{float}
\usepackage{subcaption}
\usepackage{authblk}
\usepackage{xr}

% Resolve \ref to labels defined in main.tex
\externaldocument{main}

% Custom commands (must match main.tex)
\newcommand{\model}[1]{\textsc{#1}}
\newcommand{\auc}{\text{AUC}}
\newcommand{\ci}[2]{[#1, #2]}

% Number supplementary notes/tables/figures with an S prefix
\renewcommand{\thetable}{S\arabic{table}}
\renewcommand{\thefigure}{S\arabic{figure}}
\renewcommand{\thesubsection}{S\arabic{subsection}}

\title{Supplementary Material: Stress-testing single-cell foundation models
on underrepresented platelets: a patient-level mortality benchmark}

\author[1]{Xinru Qiu}
\author[1]{Adam Godzik}
\affil[1]{Division of Biomedical Sciences, School of Medicine, University of California, Riverside, CA, USA}

\date{}

\setlength{\emergencystretch}{3em}

\begin{document}

\maketitle
\tableofcontents
\newpage

% ============================================================================
% SUPPLEMENTARY NOTES
% ============================================================================
\section{Supplementary Notes}

\subsection{Outcome-label provenance}
\label{sn:label_provenance}

Outcome labels for every fatal patient were verified against source records:
Ren (GEO outcome field), Stephenson (E-MTAB-10026 outcome attribute, 8 deaths),
Qiu GSE167363 (publication, 2 deaths), Liu (Table S6B, 4 deaths),
Schulte-Schrepping (Table S1, 3 deaths), Bernardes (publication, 1 death),
GSE279452 (series supplement, 17 deaths at 28 days), and Combes (Supplementary
Table 1 of the source publication). GSE151263 does not report per-patient
outcomes and is excluded from all mortality analyses
(Section~\ref{ssec:methods_cohort}).

\subsection{Platelet and megakaryocyte content of pretraining corpora}
\label{sn:corpus_platelet}

We measured or bounded the platelet content of each model's pretraining
corpus against public records. For \model{STATE}'s corpus the content is
bounded through its public component corpora: the scBaseCount component
(71.7M cells) contains no platelet or megakaryocyte category in its
cell-type annotation space, and its quality-control filters ($\geq$300 genes
and $\geq$500 UMIs per cell) exclude typical anucleate platelet barcodes;
Tahoe-100M consists of cancer cell lines; the Census-derived remainder
shares the Census bound below, capping \model{STATE}'s overall platelet
content below 0.05\%. For \model{UCE}'s IMA corpus, the published
51-category composition of the 2{,}969{,}114-cell representative sample
contains neither platelets nor megakaryocytes. For the three Census-derived
corpora, platelet-annotated cells are present and similarly rare:
approximately 0.120\% (\model{UCE}), 0.125\% (\model{scGPT}), and
0.120--0.125\% (\model{scPRINT-2}) of the respective releases, tens of
thousands of annotated cells. For \model{Geneformer} (Genecorpus-104M),
cell-type-resolved composition is not published in any material we could
identify and cannot be bounded. Per-corpus metrics in one unit where
available are given in Supplementary Table~\ref{tab:corpus_fractions}. This
measurement bears directly on the attribution hypothesis in
Section~\ref{ssec:discussion_property}: the three bounded models share
essentially identical platelet fractions yet span most of the between-model
performance spread, one leader (\model{STATE}) is bounded near zero, and the
other (\model{Geneformer}) cannot be bounded at all --- corpus platelet
content does not track performance wherever it can actually be checked.

\subsection{Aggregation sensitivity}
\label{sn:aggregation}

Patient-level scores were computed as the mean of a patient's out-of-fold cell
predictions. Recomputing all headline tables with median and 10\% trimmed-mean
aggregation changed no ranking and no significance call
(Supplementary Table~\ref{tab:aggregation_sens}).

\subsection{Winner-corrected primary test}
\label{sn:winner_test}

The pre-specified primary question, whether the best foundation model beats
the from-scratch baseline, was tested with a max-statistic permutation test on
the seed-ensemble out-of-fold scores: mortality labels were permuted at the
patient level (10{,}000 permutations), the maximum
foundation-model-minus-baseline AUC difference was recorded for each
permutation, and the familywise $P$ is the fraction of permutations whose
maximum reaches the observed maximum. This construction simulates the null
while preserving the inter-model score correlations (Pearson $r = 0.71$--$0.90$;
Supplementary Table~\ref{tab:score_correlation}) and the winner-selection
procedure itself; because labels are permuted on frozen scores without
retraining, it tests the slightly weaker null that the scores carry no
outcome information. The observed maximum contrast (\model{Geneformer},
$+0.064$ AUC) gave $P = 0.181$ (Supplementary
Table~\ref{tab:winner_ensemble}); a
recentered-bootstrap cross-check (subtracting the observed per-model
differences before taking the maximum in each patient resample) gave
$P = 0.073$. The pre-specified primary contrast is therefore not met. As an
exploratory comparison, a consensus predictor averaging the five models'
scores reached AUC 0.768, above the PCA-200 baseline (0.742) but below the
best single model (\model{Geneformer} 0.807): averaging over the weaker
candidates scores 0.039 AUC below the post-hoc winner on the same patients
used to identify it; this comparison is descriptive, and a prospective choice
between winner selection and ensembling would require nested evaluation.

\subsection{Variance-estimation consistency check}
\label{sn:se_check}

Pairwise AUC contrasts are reported with exact DeLong standard errors
(placement-value covariance) computed on the seed-ensemble scores. As a
consistency check, patient-level paired bootstrap standard errors agree with
the exact DeLong values for all 15 within-COVID pairs (within about 5\%;
Supplementary Table~\ref{tab:se_check}). The Hanley--McNeil formula \citep{hanley1982meaning} with the
raw-score Pearson correlation is an unpaired approximation and overestimates
the paired standard errors by 20--35\% in 8 of 15 pairs; it is not used for
any reported contrast. The winner-selection and range statistics
(Supplementary Notes~\ref{sn:winner_test} and~\ref{sn:range_test}) do not
admit a DeLong form and use permutation nulls throughout.

\subsection{Stability of the primary result}
\label{sn:fragility}

Leave-one-fatal-patient-out refits (25 in total) repeated the permutation
max-statistic test on 239 patients. Every refit returned $P > 0.05$ (range
0.141--0.288): the primary null is not driven by any single death
(Section~\ref{ssec:r1_within_covid}).

\subsection{QC-metrics-only baseline}
\label{sn:qc_baseline}

To test whether predictions could be explained by library-quality or
processing artifacts rather than biology, we trained logistic regression on
per-patient QC metrics (mean UMI counts, detected genes, percent mitochondrial
reads, and cell count) under the identical repeated-CV protocol
(Section~\ref{ssec:methods_cv}). On the COVID-only cohort (240 patients, 25
fatal) the model scored below chance (0.376 as-fit). The below-chance value is
a textbook suppression pattern, not an anomaly: the QC features are
mechanically correlated with library depth and carry outcome information only
conditional on it. Adding each patient's log-transformed cell count to the
same model recovers a correctly oriented 0.618 as-fit, and adding cohort
identity on top reaches 0.627; it is the cell count, not cohort, that flips
the sign (cohort contributes $+0.009$). Fatal patients showed higher mean UMI
counts and mitochondrial fractions than survivors in both diseases, consistent
with more activated, fragment-rich platelet transcriptomes (per-patient QC
distributions are available as machine-readable data; see Data and Code
Availability). Because a baseline race cannot separate these contributions,
the confound question is settled by the incremental-value test in
Supplementary Note~\ref{sn:qc_incremental}.

\subsection{Incremental value over QC, cell-count, and cohort covariates}
\label{sn:qc_incremental}

The decisive confound test is incremental value, not a baseline race. The
baseline covariate set (per-patient mean UMI counts, detected genes, percent
mitochondrial reads, log-transformed cell count, and cohort indicator
variables; 13 covariates) was fit with L2-regularized logistic regression
($C = 0.3$; regularization is required at 13 covariates and 25 events) under
the identical repeated-CV protocol, reaching a patient-level AUC of 0.627
as-fit on the seed-ensemble scores.

The pre-specified primary test for this question is a single nested
likelihood-ratio test, covariates-only against covariates plus the PCA-200
representation score, chosen because it involves no model selection and is
exactly the question asked: does a platelet representation add over QC, cell
count, and cohort? Because an L2-penalized statistic has no valid $\chi^2(1)$
reference, the null distribution is empirical: patient-level label
permutation (10{,}000 permutations) with both nested models refit per
permutation at the identical penalty. The observed statistic (4.477) exceeds
every permuted maximum, $P < 0.0001$ at the estimator's resolution floor
($P < 1/10{,}001$). This pre-specified test alone carries the conclusion that
mortality is learnable from platelet transcriptomes beyond QC, cell count,
and cohort.

The five foundation models are reported as exploratory, with two readings
paired so that neither is quoted alone (Supplementary
Table~\ref{tab:incremental_value}). By the same nested likelihood-ratio test
against the empirical null, every representation added significant value
(four of five at the resolution floor $P < 0.0001$; \model{UCE} weakest at
$P = 0.0022$); by paired-AUC comparison against the same baseline, five of
six increments were significant (\model{STATE} $+0.065$, \model{Geneformer}
$+0.053$, PCA-200 $+0.050$, \model{scPRINT-2} $+0.046$, \model{scGPT}
$+0.026$; exact DeLong $p = 0.0001$--$0.006$), with \model{UCE} again the
exception ($+0.016$, $p = 0.072$). The likelihood-ratio family establishes
that incremental information exists; the AUC increments quantify its
magnitude; the penalized statistic is sensitive to presence but nearly
insensitive to magnitude, so no effect size should be read from the
likelihood-ratio $P$-values themselves. Two diagnostic tables accompany the
result: the cross-validation folds are cohort-balanced (Supplementary
Table~\ref{tab:fold_balance}), and the per-cohort QC--mortality associations
are reported in Supplementary Table~\ref{tab:qc_by_cohort}.

\subsection{Each method against the count-plus-disease null}
\label{sn:method_vs_null}

Each method's seed-ensemble scores were compared with the null model
(per-patient cell count and disease label; AUC 0.554) by exact DeLong paired
tests (Supplementary Table~\ref{tab:method_vs_null}). Five of six methods
significantly outscored the null: \model{STATE}
($p = 0.0014$), \model{Geneformer} ($p = 0.0006$), \model{scPRINT-2}
($p = 0.011$), \model{scGPT} ($p = 0.023$), and PCA-200 ($p = 0.013$), all
surviving Benjamini--Hochberg correction across the six tests. \model{UCE}
did not ($+0.122$, $p = 0.157$) and is the only method failing that
correction; it is the weakest method throughout this benchmark
(Section~\ref{ssec:r1_within_covid}).

\subsection{Cohort-only baseline}
\label{sn:cohort_baseline}

Because cohort is the variable the QC diagnostic points at, we scored each
patient by their cohort's mortality rate, estimated within training folds
only with add-one (Laplace) smoothing, under the identical repeated-CV
protocol. The cohort-only model reached a patient-level AUC of
0.486 $\pm$ 0.038 (repeat-mean $\pm$ repeat-sd), statistically
indistinguishable from chance. This number must not be read as exonerating
cohort: add-one smoothing is size-dependent, so a small zero-death cohort
receives a higher smoothed rate than a large cohort with real deaths,
inverting the risk ordering by construction (per-cohort counts and smoothed
rates in Supplementary Table~\ref{tab:cohort_rates}). The estimate therefore
measures the failure of one particular cohort-risk estimator and is silent on
the confound question; the evidence on that question is the incremental-value
test with cohort identity in the covariate set
(Supplementary Note~\ref{sn:qc_incremental}).

\subsection{Patient-matched monocyte comparison}
\label{sn:monocyte_matched}

The monocyte--platelet comparison was repeated on the patient intersection
identified through the verified sample-to-patient crosswalk (Ren: sample title
base to patient; Stephenson: individual identifier; GSE167363: internal sample
map; one Ren survivor patient, S-S072, is not represented in the monocyte
cell set). Apparent direction flips under single-seed splits were traced to
split-to-split variability and disappeared under the 10-seed ensemble:
\model{STATE}, \model{scGPT}, and the PCA-200 baseline returned to the
platelet-favoring direction (or to zero difference) in all configurations,
while \model{UCE}'s monocyte-versus-platelet difference remained the closest
to zero of any method in every configuration --- the model-specific pattern
reported in Section~\ref{ssec:r4_monocyte}.

\subsection{Severe-excluded sensitivity analysis}
\label{sn:severe}

Because severe cases (ICU-discharged survivors) are clinically adjacent to
fatal outcomes and severity grading is not harmonized across source cohorts,
the primary label definition includes severe cases as survivors. The
bidirectional transfer analysis was recomputed with severe cases excluded
from the survivor class (Supplementary Table~\ref{tab:severe_sens}).
Reverse-direction point estimates shifted by $+0.011$--$+0.079$, and the
variant reverses the top two reverse-direction models (\model{Geneformer}
0.775 above \model{STATE} 0.765, with \model{UCE} 0.772 between them) ---
the basis for reporting the reverse-direction ranking as rank-indeterminate
(Section~\ref{ssec:r3_cross_disease}).

\subsection{Nested-comparison (complementarity) family}
\label{sn:nested_comparison}

A question distinct from superiority is complementarity: whether a
foundation-model score adds information beyond the baseline's score even
without ranking patients better on its own. We tested it with nested models:
a reduced model containing the PCA-200 score (in three comparator
configurations: the PCA-200 score alone; plus log-transformed cell count;
plus the full QC and cohort covariate set) against a full model adding each
foundation-model score in turn, refit inside the repeated-CV protocol with
all preprocessing fit within training folds, including per-fold feature
standardization. Any nested out-of-fold pooling harness of this kind is
required to pass a unit test before use: a known frozen score refit through
the harness must reproduce its own out-of-fold AUC within 0.02--0.04, and a
harness failing the gate produces no significance conclusions. With the gated
harness, the reduced models reached 0.710--0.724 patient-level AUC against
PCA-200's own 0.725, and no foundation model's increment survived
Benjamini--Hochberg correction ($\alpha = 0.05$) in any of the three
comparator configurations (0/5 in each; rank-by-rank values in Supplementary
Table~\ref{tab:nested_oof}); the only nominal signal was \model{STATE} at
$p = 0.025$--$0.045$ across configurations, not surviving correction. A
winner-corrected deviance permutation test on the same nested construction is
retained as an in-sample cross-check (Supplementary
Table~\ref{tab:nested_deviance}); because it does not admit an out-of-fold
form, it carries no headline weight.

\subsection{Confirmatory range-statistic test for the leaders--UCE gap}
\label{sn:range_test}

The single clearly resolvable gap in the within-COVID ranking, the two
leaders' advantage over \model{UCE} ($+0.118$ and $+0.130$), was tested
against the winner-selection null with a range statistic on the same 10{,}000
patient-level label permutations as the primary analysis: the max-minus-min
AUC across the five foundation models was recorded per permutation, and the
observed range (0.130) exceeded the null distribution
($P = 0.011$). We report this separation as the paper's single confirmatory
secondary result (Section~\ref{ssec:r1_within_covid}).

\subsection{The 5-patient GSE167363 test set}
\label{sn:gse167363_testset}

For completeness, the cross-disease transfer analysis was also run on the
5-patient GSE167363 sepsis test set (2 fatal). Forward point estimates were
substantially higher (up to 1.00 patient-level AUC: \model{STATE} 1.00,
\model{Geneformer} 1.00, \model{scGPT} 0.833, \model{UCE} 0.667,
\model{scPRINT-2} 0.667). The reverse configuration trained on the same five
patients and tested on COVID-19 (240 patients, 25 fatal) gave \model{STATE}
0.799 (95\% CI [0.718, 0.876]) with logistic regression, \model{scPRINT-2}
0.695, \model{Geneformer} 0.668, \model{scGPT} 0.661, and \model{UCE} 0.653.
Given that the pool contains five patients, these values are reported only as
a labeled small-sample reference and must not be compared with the 91-patient
and 32-patient test sets used for all reported transfer conclusions. Two further
models, TranscriptFormer \citep{pearce2025transcriptformer} and STACK, were
evaluated under a different held-out test split; their scores are not comparable to the reported transfer and
adaptation values, and they are excluded from all such comparisons
(Section~\ref{ssec:r3_cross_disease}).

\subsection{Per-patient cell-count association}
\label{sn:cellcount}

Per-patient platelet cell counts correlated only weakly with predicted scores
across models (Spearman $\rho = 0.18$--0.29 for all five foundation models and
the PCA-200 baseline; minimum 0.183 for \model{UCE}). A null model using only
cell count and disease label reached 0.554 $\pm$ 0.046 patient-level AUC within
COVID-19 (Supplementary Figure~\ref{fig:supp_cellcount} and Supplementary
Table~\ref{tab:cellcount}).

\subsection{Age structure of the GSE279452 sepsis cohort}
\label{sn:age}

The analysis set and the GEO series use three different counts, which we
state once to prevent misreading: 93 is the number of Adult Sepsis candidate
donors in GSE279452; 92 of them matched per-sample age records directly (one
donor, \texttt{Skin-S152}, matches only after correcting a sample-title
spelling inconsistency in the GEO series, where the sample is titled
\texttt{Skin-152} while its same-site siblings carry the \texttt{Skin-S}
prefix); and 91 patients enter the analysis set (93 minus two donors lacking
the 28-day outcome label), of whom 90 have direct age records. All known ages
are adult (median 67 years, range 19--93; the two youngest donors are 19
years old). The full GSE279452 series additionally contains 26 pediatric
samples (17 sepsis, 9 healthy or ICU control; ages 0.1--15 years); none of
them entered the candidate set, so the pediatric/adult mixture and the
site--age correlation present in the full series do not enter any reported
analysis. Fatal outcomes stratify by age as expected:
0 of 6 patients below 40, 7 of 30 patients aged 40--60, and 10 of 54
patients over 60 with a known 28-day outcome (2 further donors in this stratum
lack the 28-day label; 17 fatal in total). The cohort is single-site and single-platform
(10x), which is reported as a cohort-level confound for the cross-disease
transfer analyses (Section~\ref{ssec:limitations}; Supplementary
Figure~\ref{fig:supp_age}).

\subsection{Masking-based gene attribution and cross-model GSEA}
\label{sn:attribution_gsea}

The masking-attribution pipeline (Section~\ref{ssec:methods_masking}) produced
one ranked per-gene list per model (5{,}038 genes for \model{Geneformer},
5{,}000 for \model{STATE}, 4{,}987 for \model{scGPT}, 4{,}850 for \model{UCE},
4{,}850 for \model{scPRINT-2}, each passing a minimum-occurrence filter of 30
cells), of which 3{,}917 genes are ranked by all five models and form the
common comparison set used for every cross-model statistic in
Section~\ref{ssec:r6_attribution_gsea}. Pairwise Spearman rank correlations
of the full lists range from $-0.03$ to $0.40$; 531 of the 3{,}917 common
genes (365 positive, 166 negative) carry the same sign in all five models,
exceeding the $\sim$272 expected under the models' independent marginal sign
rates. The per-model GSEA tables (Hallmark 2020, KEGG 2021, Reactome 2022;
FDR $q < 0.25$) contain 6/23/2 significant terms for \model{Geneformer},
5/7/2 for \model{STATE}, 28/63/257 for \model{scGPT}, 23/70/82 for
\model{UCE}, and 0/0/0 for \model{scPRINT-2}; no term is significant in four
or more models. The joint-occlusion dose control masked the top-$N$ attributed
genes jointly ($N \in \{10, 30, 128\}$, up to half of each cell's candidate
sentence) against matched random-$N$ sets (three draws per cell, $n = 10{,}736$
cells): joint deltas were indistinguishable from zero and from random-$N$
occlusion across all tiers and polarities (paired tests, $P = 0.21$--$0.93$),
while the additive prediction from single-gene attributions was 55--750 times
larger. As stated in Section~\ref{ssec:r6_attribution_gsea}, the candidate
gene universes differed across models by construction (the vocabulary
intersection of the models built up to each run: 11{,}984 genes for
\model{Geneformer} and \model{STATE}, 11{,}977 after intersecting
\model{scGPT}'s vocabulary, smaller still after \model{UCE}'s and
\model{scPRINT-2}'s), and \model{UCE}'s attributions used the one-line
float-to-boolean mask correction described in
Section~\ref{ssec:methods_masking}. Full ranked lists, GSEA tables, the
cross-model comparison tables, and the occlusion-control per-cell results are
available as machine-readable data (see Data and Code Availability).


% ============================================================================
% SUPPLEMENTARY TABLES
% ============================================================================
\section{Supplementary Tables}

% tab:corpus_fractions — one-unit corpus metrics for all five models (main text)
\begin{table}[t]
\centering
\caption{Pretraining-corpus composition in one unit per model: the fraction of
pretraining cells from blood/bone marrow, from inflammatory or infectious
conditions, and annotated as platelets, as documented in each corpus's own
records. ``Not reported'' means no cell-type- or condition-resolved
composition is published for that corpus. Checkpoint names are pinned because
corpus claims differ across versions.}
\label{tab:corpus_fractions}
\small
\resizebox{\textwidth}{!}{%
\begin{tabular}{llllll}
\toprule
\textbf{Model (checkpoint)} & \textbf{Corpus} & \textbf{Cells} & \textbf{Blood / marrow} & \textbf{Inflamm./infect.} & \textbf{Platelet} \\
\midrule
STATE (SE-600M) & scBaseCount + CxG + Tahoe-100M & 167M & largest single share of scBaseCount (43\% of corpus), exact \% not reported; Tahoe part is cancer cell lines & 30.7\% cancer within scBaseCount; rest not reported & structurally absent in scBaseCount (no platelet category); full corpus not reported \\
Geneformer (V2-316M) & Genecorpus-104M (non-cancer) & 104M & $\leq$25\% by construction (per-tissue cap); 49 of 561 source datasets are bone marrow & disease conditions included (only high-mutational-burden cells excluded); \% not reported & not reported \\
scPRINT-2 (small-v2) & CxG + Tahoe-100M + scBaseCount & 350M & not reported (scBaseCount part shares scBaseCount's blood/marrow skew) & not reported & 0.120--0.125\% (Census-derived part) \\
scGPT (whole-human) & CELLxGENE, non-disease & 33M & not reported (51 organs incl. blood and bone marrow) & 0\% by design (non-disease only) & 0.125\% \\
UCE (650M) & IMA (CxG + 28 datasets) & 36M & 6.6\% (SI Table 19 representative sample); brain $\approx$25\% & not reported (mostly reference atlases) & 0.120\% (Census-derived part); published 51-category list: none \\
\bottomrule
\end{tabular}%
}
\end{table}
% tab:property_multiplicity — candidate attributes and whether they separate the observed ranking (supplementary)
\begin{table}[t]
\centering
\caption{Candidate model attributes and whether each separates the observed
ranking, taking the clearly resolvable gap to be \model{Geneformer}/\model{STATE}
(top two, statistically inseparable) versus \model{scGPT}/\model{UCE} (bottom
two). \model{scPRINT-2} sits between the groups and is shown in every row
because it breaks several otherwise-clean separations. The composite property
discussed in Section~\ref{ssec:discussion_property} is one of the few
attributes that fits, and the only one with a direct measurement, which
falsifies it where it is measurable
(Supplementary Note~\ref{sn:corpus_platelet}).}
\label{tab:property_multiplicity}
\small
\resizebox{\textwidth}{!}{%
\begin{tabular}{llll}
\toprule
\textbf{Candidate attribute} & \textbf{Separates top-2 from bottom-2?} & \textbf{Direction} & \textbf{Contradiction} \\
\midrule
Corpus blood / bone-marrow fraction & \textbf{yes} & leaders blood/marrow-skewed; \model{scGPT} mixed, \model{UCE} 6.6\% & \model{scPRINT-2} shares the scBaseCount skew yet ranks third; the platelet-specific fraction is falsified where measurable (Supplementary Note~\ref{sn:corpus_platelet}) \\
Corpus disease-condition fraction & partial & \model{scGPT} 0\% (fifth) & \model{UCE} mixed-corpus yet last \\
Objective: rank-ordering of genes & no & \model{Geneformer} only & separates \model{Geneformer} alone, not the pair \\
Objective includes population-level term & no & \model{STATE} only & separates \model{STATE} alone \\
Objective: per-gene value reconstruction & no & \model{STATE}, \model{scGPT}, \model{UCE} share it & groups a leader with both bottom models \\
Tokenization: protein-embedding genes & no & \model{STATE}, \model{UCE}, \model{scPRINT-2} & mixes leader, bottom, and middle \\
Supervised label heads in pretraining & no & \model{scPRINT-2} only & middle rank, breaks the identity story \\
Corpus size (cells) & no & \model{scPRINT-2} largest (350M), mid-ranked & --- \\
Parameter count & no & \model{UCE} largest (675M), last & --- \\
Embedding dimension & no & \model{STATE} 2{,}058 first; \model{UCE} 1{,}280 last & --- \\
Species coverage & no & \model{UCE}, \model{scPRINT-2} multi-species & one last, one mid \\
Release recency & no & no monotone relation & --- \\
\bottomrule
\end{tabular}%
}
\end{table}

\begin{table}[t]
\centering
\caption{Incremental value over the QC + cell count + cohort covariate
baseline (patient-level AUC 0.627 as-fit; 240 patients, 25 fatal). The
PCA-200 row is the pre-specified primary for this question (nested
likelihood-ratio test against a 10{,}000-permutation empirical null, chosen
because it involves no model selection); the five foundation models are
exploratory, with both readings paired. The nested-LR family establishes
that incremental information exists; the $\Delta$AUC column quantifies its
magnitude.}
\label{tab:incremental_value}
\small
\begin{tabular}{lccccc}
\toprule
\textbf{Model} & \textbf{Augmented AUC} & \textbf{$\Delta$AUC} & \textbf{DeLong SE} & \textbf{DeLong $p$} & \textbf{Nested-LR $P$} \\
\midrule
\model{STATE} & 0.692 & $+0.065$ & 0.0171 & 0.0002 & $<0.0001$ \\
\model{Geneformer} & 0.680 & $+0.053$ & 0.0140 & 0.0001 & $<0.0001$ \\
PCA-200 (pre-specified primary) & 0.677 & $+0.050$ & 0.0172 & 0.0034 & $<0.0001$ \\
\model{scPRINT-2} & 0.673 & $+0.046$ & 0.0167 & 0.0054 & $<0.0001$ \\
\model{scGPT} & 0.653 & $+0.026$ & 0.0097 & 0.0063 & $<0.0001$ \\
\model{UCE} & 0.643 & $+0.016$ & 0.0089 & 0.072 & 0.0022 \\
\bottomrule
\end{tabular}
\end{table}

\begin{table}[t]
\centering
\caption{Each method's seed-ensemble scores against the count-plus-disease
null model (AUC 0.554; exact DeLong paired tests, 240 patients, 25 fatal).
Five of six methods survive Benjamini--Hochberg correction across the six
tests; \model{UCE} is the only one that does not.}
\label{tab:method_vs_null}
\small
\begin{tabular}{lcccc}
\toprule
\textbf{Method} & \textbf{AUC} & \textbf{$\Delta$ vs null} & \textbf{DeLong SE} & \textbf{DeLong $p$} \\
\midrule
\model{Geneformer} & 0.807 & $+0.252$ & 0.0738 & 0.0006 \\
\model{STATE} & 0.795 & $+0.240$ & 0.0749 & 0.0014 \\
\model{scPRINT-2} & 0.756 & $+0.202$ & 0.0796 & 0.011 \\
PCA-200 & 0.742 & $+0.188$ & 0.0753 & 0.013 \\
\model{scGPT} & 0.739 & $+0.185$ & 0.0814 & 0.023 \\
\model{UCE} & 0.677 & $+0.122$ & 0.0863 & 0.157 \\
\bottomrule
\end{tabular}
\end{table}

\begin{table}[t]
\centering
\caption{Per-cohort mortality rates and the add-one (Laplace) smoothed rate
used by the cohort-only baseline. Smoothing is size-dependent: small
zero-death cohorts (Wilk, Lee, Arunachalam, Combes) receive smoothed rates
above Schulte-Schrepping's and Bernardes's raw rates, inverting the risk
ordering by construction; the cohort-only AUC of 0.486 therefore measures
the failure of this estimator, not the absence of cohort signal.}
\label{tab:cohort_rates}
\small
\begin{tabular}{lrrcc}
\toprule
\textbf{Cohort} & \textbf{Patients} & \textbf{Fatal} & \textbf{Raw rate} & \textbf{Laplace rate} \\
\midrule
Stephenson & 90 & 8 & 0.089 & 0.098 \\
Ren & 54 & 9 & 0.167 & 0.179 \\
Liu & 31 & 4 & 0.129 & 0.152 \\
Schulte-Schrepping & 18 & 3 & 0.167 & 0.200 \\
Bernardes & 13 & 1 & 0.077 & 0.133 \\
Combes & 13 & 0 & 0.000 & 0.067 \\
Lee & 8 & 0 & 0.000 & 0.100 \\
Arunachalam & 7 & 0 & 0.000 & 0.111 \\
Wilk & 6 & 0 & 0.000 & 0.125 \\
\bottomrule
\end{tabular}
\end{table}

\begin{table}[t]
\centering
\caption{Cohort balance across cross-validation folds (seed 42, one
representative repeat): every test fold contains 48 patients from all nine
cohorts, with 4--6 fatal patients. Per-fold per-cohort counts are available
as machine-readable data (see Data and Code Availability).}
\label{tab:fold_balance}
\small
\begin{tabular}{cccc}
\toprule
\textbf{Fold} & \textbf{Test patients} & \textbf{Test fatal} & \textbf{Cohorts in test} \\
\midrule
1 & 48 & 4 & 9 \\
2 & 48 & 4 & 9 \\
3 & 48 & 6 & 9 \\
4 & 48 & 5 & 9 \\
5 & 48 & 6 & 9 \\
\bottomrule
\end{tabular}
\end{table}

\begin{table}[t]
\centering
\caption{Per-cohort association between QC metrics (median UMI counts) and
mortality (Mann--Whitney $U$, uncorrected). Cohorts with zero fatal patients
cannot be tested. One of six testable cohorts reaches nominal $p = 0.031$;
none survives multiple-comparison correction.}
\label{tab:qc_by_cohort}
\small
\begin{tabular}{lccccc}
\toprule
\textbf{Cohort} & \textbf{Patients} & \textbf{Fatal} & \textbf{Median UMI (fatal)} & \textbf{Median UMI (surv.)} & \textbf{MWU $p$} \\
\midrule
Stephenson & 90 & 8 & 778 & 980 & 0.288 \\
Ren & 54 & 9 & 1{,}682 & 1{,}684 & 0.790 \\
Liu & 31 & 4 & 1{,}226 & 693 & 0.031 \\
Schulte-Schrepping & 18 & 3 & 3{,}402 & 1{,}299 & 0.100 \\
Bernardes & 13 & 1 & 533 & 550 & 1.000 \\
Wilk, Lee, Arunachalam, Combes & 6--13 & 0 & \multicolumn{3}{c}{insufficient variation for within-cohort test} \\
\bottomrule
\end{tabular}
\end{table}

\begin{table}[t]
\centering
\caption{Nested-comparison (complementarity) family: out-of-fold
$\Delta$AUC (exact DeLong $p$) of each foundation-model score over three
reduced-model comparators, refit in the repeated-CV protocol with per-fold
standardization (the gated harness; reduced-model OOF AUCs now sit at
0.710--0.724 against PCA-200's own 0.725). No increment survives
Benjamini--Hochberg correction ($\alpha = 0.05$) in any configuration (0/5
each); \model{STATE}'s nominal $p = 0.025$--$0.045$ is the only signal
below 0.05 and does not survive correction.}
\label{tab:nested_oof}
\small
\resizebox{\textwidth}{!}{%
\begin{tabular}{lccccc}
\toprule
\textbf{Reduced model (OOF AUC)} & \textbf{\model{STATE} $\Delta$ ($p$)} & \textbf{\model{Geneformer} $\Delta$ ($p$)} & \textbf{\model{scPRINT-2} $\Delta$ ($p$)} & \textbf{\model{scGPT} $\Delta$ ($p$)} & \textbf{\model{UCE} $\Delta$ ($p$)} \\
\midrule
QC + count + cohort + PCA-200 (0.710) & $+0.024$ (0.031) & $+0.033$ (0.083) & $-0.002$ (0.891) & $+0.000$ (0.967) & $-0.000$ (0.906) \\
PCA-200 alone (0.724) & $+0.052$ (0.045) & $+0.049$ (0.056) & $+0.019$ (0.402) & $+0.010$ (0.496) & $-0.007$ (0.306) \\
PCA-200 + log cell count (0.713) & $+0.060$ (0.025) & $+0.056$ (0.045) & $+0.028$ (0.264) & $+0.013$ (0.361) & $-0.008$ (0.271) \\
\bottomrule
\end{tabular}%
}
\end{table}

\begin{table}[t]
\centering
\caption{In-sample deviance cross-check on the nested construction
(winner-corrected permutation null over 10{,}000 label permutations: median
0.188, 95th percentile 0.790, 99th percentile 1.230, maximum 3.819).
Retained as an in-sample cross-check only; because this construction does
not admit an out-of-fold form, it carries no headline weight
(Supplementary Note~\ref{sn:nested_comparison}).}
\label{tab:nested_deviance}
\small
\begin{tabular}{lcc}
\toprule
\textbf{Model} & \textbf{Deviance statistic} & \textbf{Position vs null} \\
\midrule
\model{STATE} & 4.328 & exceeds null maximum \\
\model{Geneformer} & 3.223 & exceeds 99th percentile \\
\model{scPRINT-2} & 3.184 & exceeds 99th percentile \\
\model{scGPT} & 1.537 & exceeds 99th percentile \\
\model{UCE} & 0.875 & exceeds 95th percentile \\
\bottomrule
\end{tabular}
\end{table}

\begin{table}[t]
\centering
\caption{Pearson correlations between the six methods' seed-ensemble
patient-level scores (within-COVID). High inter-model correlations
($r = 0.71$--$0.90$) are what make the winner-correction in the primary
test necessary (Supplementary Note~\ref{sn:winner_test}).}
\label{tab:score_correlation}
\small
\begin{tabular}{lcccccc}
\toprule
 & \textbf{\model{STATE}} & \textbf{\model{UCE}} & \textbf{\model{Geneformer}} & \textbf{\model{scPRINT-2}} & \textbf{\model{scGPT}} & \textbf{PCA-200} \\
\midrule
\model{STATE} & 1.000 & 0.794 & 0.904 & 0.897 & 0.816 & 0.859 \\
\model{UCE} & 0.794 & 1.000 & 0.734 & 0.755 & 0.843 & 0.707 \\
\model{Geneformer} & 0.904 & 0.734 & 1.000 & 0.877 & 0.781 & 0.830 \\
\model{scPRINT-2} & 0.897 & 0.755 & 0.877 & 1.000 & 0.759 & 0.784 \\
\model{scGPT} & 0.816 & 0.843 & 0.781 & 0.759 & 1.000 & 0.776 \\
PCA-200 & 0.859 & 0.707 & 0.830 & 0.784 & 0.776 & 1.000 \\
\bottomrule
\end{tabular}
\end{table}

\begin{table}[t]
\centering
\caption{Per-model advantage over the PCA-200 baseline on the seed-ensemble
scores with patient-level bootstrap intervals (10{,}000 resamples) and
nominal uncorrected $p$-values. The pre-specified primary evaluates only the
maximum of the five foundation-model contrasts (\model{Geneformer},
$+0.064$) against the permutation max-statistic null: familywise
$P = 0.181$ (Supplementary Note~\ref{sn:winner_test}).}
\label{tab:winner_ensemble}
\small
\begin{tabular}{lcccc}
\toprule
\textbf{Model} & \textbf{AUC} & \textbf{$\Delta$ vs PCA-200} & \textbf{Bootstrap 95\% CI} & \textbf{Uncorrected $p$} \\
\midrule
\model{Geneformer} & 0.807 & $+0.064$ & [0.007, 0.127] & 0.026 \\
\model{STATE} & 0.795 & $+0.052$ & [0.004, 0.107] & 0.033 \\
\model{scPRINT-2} & 0.756 & $+0.014$ & [$-0.044$, 0.078] & 0.675 \\
\model{scGPT} & 0.739 & $-0.003$ & [$-0.057$, 0.053] & 0.894 \\
\model{UCE} & 0.677 & $-0.066$ & [$-0.140$, 0.009] & 0.087 \\
\bottomrule
\end{tabular}
\end{table}

\begin{table}[t]
\centering
\caption{Variance-estimation consistency: patient-level paired bootstrap
versus exact DeLong standard errors for all 15 within-COVID pairwise
contrasts (seed-ensemble scores). The two estimators agree within about 5\%
throughout (Supplementary Note~\ref{sn:se_check}).}
\label{tab:se_check}
\small
\begin{tabular}{lcccc}
\toprule
\textbf{Pair} & \textbf{$\Delta$AUC} & \textbf{Bootstrap SE} & \textbf{Exact DeLong SE} & \textbf{DeLong $p$} \\
\midrule
\model{STATE} -- \model{UCE} & 0.118 & 0.0279 & 0.0278 & $<$0.0001 \\
\model{STATE} -- \model{Geneformer} & $-0.012$ & 0.0220 & 0.0217 & 0.577 \\
\model{STATE} -- \model{scPRINT-2} & 0.039 & 0.0194 & 0.0190 & 0.043 \\
\model{STATE} -- \model{scGPT} & 0.055 & 0.0229 & 0.0225 & 0.014 \\
\model{STATE} -- PCA-200 & 0.052 & 0.0266 & 0.0265 & 0.048 \\
\model{UCE} -- \model{Geneformer} & $-0.130$ & 0.0267 & 0.0262 & $<$0.0001 \\
\model{UCE} -- \model{scPRINT-2} & $-0.080$ & 0.0295 & 0.0293 & 0.0066 \\
\model{UCE} -- \model{scGPT} & $-0.063$ & 0.0254 & 0.0250 & 0.012 \\
\model{UCE} -- PCA-200 & $-0.066$ & 0.0383 & 0.0377 & 0.081 \\
\model{Geneformer} -- \model{scPRINT-2} & 0.051 & 0.0265 & 0.0259 & 0.051 \\
\model{Geneformer} -- \model{scGPT} & 0.068 & 0.0262 & 0.0256 & 0.008 \\
\model{Geneformer} -- PCA-200 & 0.064 & 0.0309 & 0.0307 & 0.036 \\
\model{scPRINT-2} -- \model{scGPT} & 0.017 & 0.0234 & 0.0228 & 0.458 \\
\model{scPRINT-2} -- PCA-200 & 0.014 & 0.0311 & 0.0308 & 0.655 \\
\model{scGPT} -- PCA-200 & $-0.003$ & 0.0278 & 0.0277 & 0.909 \\
\bottomrule
\end{tabular}
\end{table}

\begin{table}[t]
\centering
\caption{Severe-excluded sensitivity analysis, reverse transfer direction
(sepsis$\rightarrow$COVID, logistic regression; patient-level AUC). Point
estimates shift by $+0.011$--$+0.079$ and the top two models reverse, the
basis for reporting the reverse-direction ranking as rank-indeterminate
(Section~\ref{ssec:r3_cross_disease}).}
\label{tab:severe_sens}
\small
\begin{tabular}{lccc}
\toprule
\textbf{Model} & \textbf{Primary (severe as survived)} & \textbf{Severe excluded} & \textbf{Shift} \\
\midrule
\model{STATE} & 0.721 & 0.765 & $+0.044$ \\
\model{Geneformer} & 0.712 & 0.775 & $+0.063$ \\
\model{UCE} & 0.693 & 0.772 & $+0.079$ \\
\model{scGPT} & 0.642 & 0.678 & $+0.036$ \\
\model{scPRINT-2} & 0.619 & 0.630 & $+0.011$ \\
\bottomrule
\end{tabular}
\end{table}

\begin{table}[t]
\centering
\caption{Aggregation sensitivity: patient-level AUC (repeat-mean $\pm$
repeat-sd) with mean versus median aggregation of each patient's out-of-fold
cell predictions (logistic regression). The 10\% trimmed-mean variant is
intermediate. No ranking and no significance call changes across
aggregation rules (Supplementary Note~\ref{sn:aggregation}).}
\label{tab:aggregation_sens}
\small
\begin{tabular}{lcc}
\toprule
\textbf{Method} & \textbf{Mean (primary)} & \textbf{Median} \\
\midrule
\model{Geneformer} & 0.786 $\pm$ 0.026 & 0.765 $\pm$ 0.029 \\
\model{STATE} & 0.778 $\pm$ 0.033 & 0.779 $\pm$ 0.036 \\
\model{scPRINT-2} & 0.740 $\pm$ 0.027 & 0.732 $\pm$ 0.029 \\
PCA-200 & 0.725 $\pm$ 0.033 & 0.696 $\pm$ 0.036 \\
\model{scGPT} & 0.720 $\pm$ 0.047 & 0.706 $\pm$ 0.050 \\
\model{UCE} & 0.661 $\pm$ 0.046 & 0.645 $\pm$ 0.048 \\
\bottomrule
\end{tabular}
\end{table}

\begin{table}[t]
\centering
\caption{Leave-one-cohort-out folds (patient-level AUC on the held-out
cohort). Folds are computable only for cohorts containing at least one
fatal and one survivor patient; all-survivor cohorts (Wilk, Lee,
Arunachalam, Combes) contribute to training throughout and are never scored
out-of-fold. The pooled estimate over concatenated held-out predictions
(206 patients, 25 fatal) is reported in
Section~\ref{ssec:r5_calibration}.}
\label{tab:lobo_folds}
\small
\begin{tabular}{llccc}
\toprule
\textbf{Model} & \textbf{Held-out cohort} & \textbf{AUC} & \textbf{Patients} & \textbf{Fatal} \\
\midrule
\model{STATE} & Schulte-Schrepping & 0.800 & 18 & 3 \\
 & Ren & 0.909 & 54 & 9 \\
 & Liu & 0.926 & 31 & 4 \\
 & Stephenson & 0.662 & 90 & 8 \\
 & Bernardes & 1.000 & 13 & 1 \\
\midrule
\model{UCE} & Schulte-Schrepping & 0.933 & 18 & 3 \\
 & Ren & 0.770 & 54 & 9 \\
 & Liu & 0.963 & 31 & 4 \\
 & Stephenson & 0.582 & 90 & 8 \\
 & Bernardes & 0.833 & 13 & 1 \\
\midrule
\model{Geneformer} & Schulte-Schrepping & 0.867 & 18 & 3 \\
 & Ren & 0.891 & 54 & 9 \\
 & Liu & 0.954 & 31 & 4 \\
 & Stephenson & 0.634 & 90 & 8 \\
 & Bernardes & 1.000 & 13 & 1 \\
\midrule
\model{scPRINT-2} & Schulte-Schrepping & 0.911 & 18 & 3 \\
 & Ren & 0.842 & 54 & 9 \\
 & Liu & 0.963 & 31 & 4 \\
 & Stephenson & 0.659 & 90 & 8 \\
 & Bernardes & 1.000 & 13 & 1 \\
\midrule
\model{scGPT} & Schulte-Schrepping & 0.733 & 18 & 3 \\
 & Ren & 0.770 & 54 & 9 \\
 & Liu & 0.944 & 31 & 4 \\
 & Stephenson & 0.613 & 90 & 8 \\
 & Bernardes & 0.833 & 13 & 1 \\
\bottomrule
\end{tabular}
\end{table}

\begin{table}[t]
\centering
\caption{Per-patient cell-count association with predicted scores and with
mortality (within-COVID, 240 patients; count terciles have 80/83/77
patients and 7/9/9 fatal). Scores correlate only weakly with cell count
(Spearman $\rho = 0.18$--$0.29$), and the count-only null model reaches
0.554 (Supplementary Note~\ref{sn:cellcount}).}
\label{tab:cellcount}
\small
\begin{tabular}{lccccc}
\toprule
\textbf{Method} & \textbf{Spearman $\rho$ (score, count)} & \textbf{$p$} & \textbf{AUC, low-count third} & \textbf{AUC, mid third} & \textbf{AUC, high third} \\
\midrule
\model{STATE} & 0.226 & 0.0004 & 0.759 & 0.811 & 0.543 \\
\model{UCE} & 0.183 & 0.004 & 0.687 & 0.716 & 0.397 \\
\model{Geneformer} & 0.250 & $<$0.0001 & 0.853 & 0.825 & 0.659 \\
\model{scPRINT-2} & 0.271 & $<$0.0001 & 0.763 & 0.754 & 0.621 \\
\model{scGPT} & 0.217 & 0.0007 & 0.748 & 0.710 & 0.405 \\
PCA-200 & 0.274 & $<$0.0001 & 0.732 & 0.728 & 0.489 \\
\bottomrule
\end{tabular}
\end{table}


% ============================================================================
% SUPPLEMENTARY FIGURES
% ============================================================================
\section{Supplementary Figures}

\begin{figure}[t]
\centering
\includegraphics[width=0.85\textwidth]{figures/v2/FigureS1_cellcount.pdf}
\caption{\textbf{Per-patient cell-count association.} Predicted
mortality-classifier scores versus per-patient platelet cell counts for each
method (within-COVID). Scores correlate only weakly with cell count
(Spearman $\rho = 0.18$--$0.29$; Supplementary
Note~\ref{sn:cellcount} and Table~\ref{tab:cellcount}), so mortality is not
readable from cell counts alone (count-only null AUC 0.554).}
\label{fig:supp_cellcount}
\end{figure}

\begin{figure}[t]
\centering
\includegraphics[width=0.95\textwidth]{figures/v2/FigureS2_pairwise_heatmap.pdf}
\caption{\textbf{Within-COVID pairwise contrasts.} Full 6$\times$6 matrix of
pairwise $\Delta$AUC (exact DeLong $p$) between the five foundation models
and the PCA-200 baseline on the seed-ensemble scores; cells surviving
Benjamini--Hochberg correction (family of 15) are marked. All six
significant contrasts have \model{UCE} or \model{scGPT} on the losing side,
and the two leaders-versus-\model{UCE} gaps are the largest
(Supplementary Table~\ref{tab:se_check}).}
\label{fig:supp_pairwise}
\end{figure}

\begin{figure}[t]
\centering
\includegraphics[width=0.85\textwidth]{figures/v2/FigureS_dropout.pdf}
\caption{\textbf{Embedding-dimension dropout robustness.} Patient-level AUC
as a function of randomly zeroed embedding dimensions (0--70\%), per model
(mortality endpoint, legacy analysis pool; mean $\pm$ 1 SD over five
patient-grouped folds). \model{UCE}'s curve is flat ($\Delta = +0.003$ at
70\% zeroing), indicating a highly distributed representation; every other
model degrades by 0.010--0.025 AUC.}
\label{fig:supp_dropout}
\end{figure}

\begin{figure}[t]
\centering
\includegraphics[width=0.95\textwidth]{figures/v2/FigureS_LOBO.pdf}
\caption{\textbf{Leave-one-cohort-out generalization.} Per-fold
patient-level AUC on each held-out cohort (computable only for cohorts
containing at least one fatal and one survivor patient; per-fold counts in
Supplementary Table~\ref{tab:lobo_folds}). The top three models sit within
0.005 AUC of each other on the pooled concatenated estimate (206 patients,
25 fatal; Section~\ref{ssec:r5_calibration}).}
\label{fig:supp_lobo}
\end{figure}

\begin{figure}[t]
\centering
\includegraphics[width=0.85\textwidth]{figures/v2/FigureS_Age_gse279452.pdf}
\caption{\textbf{Age structure of the GSE279452 sepsis cohort.} All 90
analysis-set patients with age records are adults (median 67 years, range
19--93); the full GEO series additionally contains 26 pediatric samples,
none of which entered the candidate or analysis sets (Supplementary
Note~\ref{sn:age}).}
\label{fig:supp_age}
\end{figure}


% ============================================================================
% REFERENCES CITED IN THE SUPPLEMENT
% ============================================================================
\begin{thebibliography}{9}

\bibitem[Hanley \& McNeil(1982)]{hanley1982meaning}
Hanley, J. A., \& McNeil, B. J. (1982).
\newblock The meaning and use of the area under a receiver operating characteristic (ROC) curve.
\newblock Radiology, 143(1), 29--36. doi:10.1148/radiology.143.1.7063747.

\bibitem[Pearce et al.(2026)]{pearce2025transcriptformer}
Pearce, J. D., Simmonds, S. E., Mahmoudabadi, G., et al. (2026).
\newblock TranscriptFormer: a generative cell atlas across 1.5 billion years of evolution.
\newblock Science, 393(6806), aec8514. doi:10.1126/science.aec8514.

\end{thebibliography}

\end{document}
